\chapter{Derivational Morphology}\label{ch:derivational}

% ============================================================
% STATUS: SCAFFOLD — port and revise from
% archive/2020/chapters/06-derivational-morphology.tex.
% Key review points:
%   - Nominal derivation: update case terminology in examples
%   - Agent nominals: relation to the new antipassive stem (-kou)
%   - Abstract nominals: interaction with the new case system
%   - Verbal derivation: how does it interact with the antipassive /
%     applicative / causative slot? Are there derivational prefixes that
%     interact with valency?
%   - Borrowing: update to include recent borrowings and assimilation
%     patterns under the new stem-class system (the German -irná verbs)
% ============================================================

\section{Nominal derivation}

\subsection{Diminutives and augmentatives}

% [Port from 2020 archive]

\subsection{Agent nouns: \ird{-(h)ár} and \ird{-ist}}

Agent nouns denoting the habitual or professional performer of an action are formed with the suffixes \ird{-(h)ár} and \ird{-ist} (the latter a Slavic borrowing, now productive):

% [Examples: fisherman, teacher, activist, etc.]
% Note: the agent nominal -kou (from the antipassive stem) is a verbal
% derivation documented in §\,\ref{sec:nominalized}; the -ár/-ist forms
% are purely nominal derivation.

\subsection{Place nouns: \ird{-(i)mašt}}

% [Port from 2020 archive]

\subsection{Abstract nouns: \ird{-(i)žnám}}

% [Port from 2020 archive]

\section{Verbal derivation}

% ============================================================
% STATUS: STUB — the interaction of valency operations (antipassive,
% applicative, causative) with derivational morphology needs to be
% clarified. The question is whether some "derivational" operations
% in the 2020 archive are now better analyzed as instances of the
% applicative or causative, or whether they are genuinely separate.
% ============================================================

\section{Compounding}

% [Port from 2020 archive; update examples]

\section{Borrowing and assimilation patterns}

% The 2020 archive described:
%   - Phonological assimilation of loanwords (stress, consonant substitution)
%   - German loanwords in -irná (treated as a separate verbal class:
%     thematic consonant = -r, active root in -irž-)
%   - Recent borrowings (English, international) showing less assimilation
% These patterns remain valid in the 2026 framework; update terminology only.

